\section*{Hoofdstuk \ref{ch:b3:origin-of-elements} --- Waar de elementen vandaan komen}
% Values from the Book 5 scratch calculator calc/ch30.py.

\begin{solution}{exo:b3:origin-of-elements:1}
(a)~$\mathrm{Si/Fe} = 10^{7.51 - 7.46} = 10^{0.05} = 1.12$ en
$\mathrm{O/Fe} = 10^{8.69 - 7.46} = 10^{1.23} = 17.0$: ongeveer één siliciumatoom
en zeventien zuurstofatomen per ijzeratoom. (b)~$\mathrm{Fe/Au} =
10^{7.46 - 0.91} = 10^{6.55} = \num{3.55e6}$ en $\mathrm{Pb/Au} =
10^{1.95 - 0.91} = 10^{1.04} = 11.0$: voorbij de ijzerpiek dalen de abundanties
met een factor van miljoenen, maar lood, op zijn eigen piek, is elf keer zo
overvloedig als goud.
\end{solution}

\begin{solution}{exo:b3:origin-of-elements:2}
(a)~$10^{7.46 - 12} = \num{2.884e-5}$, dat is \num{28.8} ijzeratomen per
miljoen waterstofatomen. (b)~Honderd keer minder: \num{0.288} per
miljoen, één ijzeratoom per $1/\num{2.884e-7} = \num{3.47e6}$
waterstofatomen. (c)~$[\mathrm{Fe/H}] = \log_{10}(\num{e-7}/\num{2.884e-5}) =
-2.46$.
\end{solution}

\begin{solution}{exo:b3:origin-of-elements:3}
$^{124}\mathrm{Xe}$ en $^{126}\mathrm{Xe}$ liggen aan de protonrijke kant,
buiten het s-pad, en zijn afgeschermd door de stabiele $^{124}\mathrm{Te}$ en
$^{126}\mathrm{Te}$: \omterm{def:b3:origin-of-elements:p-process}{p-kernen}. $^{128}\mathrm{Xe}$ en $^{130}\mathrm{Xe}$
liggen op het s-pad maar zijn voor het \omterm{def:b3:origin-of-elements:s-process}{r-proces} afgeschermd door de stabiele
$^{128}\mathrm{Te}$ en $^{130}\mathrm{Te}$: alleen het \omterm{def:b3:origin-of-elements:s-process}{s-proces}. $^{129}\mathrm{Xe}$,
$^{131}\mathrm{Xe}$ en $^{132}\mathrm{Xe}$ liggen op het s-pad en hebben geen stabiele
isobaar met lager $Z$ (jodium heeft er alleen 127): beide processen.
$^{134}\mathrm{Xe}$ en $^{136}\mathrm{Xe}$ liggen voorbij
$^{133}\mathrm{Xe}$, waar het s-pad xenon verlaat, en zijn niet
afgeschermd: alleen door het \omterm{def:b3:origin-of-elements:s-process}{r-proces}.
\end{solution}

\begin{solution}{exo:b3:origin-of-elements:4}
(a)~$t = d/c = 40\times\qty{3.086e22}{m}/\qty{2.998e8}{m.s^{-1}} =
\qty{4.12e15}{s}$, ongeveer \num{130} miljoen jaar. (b)~$\Delta v/c \le
1.7/\num{4.12e15} = \num{4.1e-16}$. Dit is een bovengrens, omdat een deel of
de hele vertraging kan voortkomen uit de uitzending van de gammastraling na de
samensmelting.
\end{solution}

\begin{solution}{exo:b3:origin-of-elements:5}
(a)~$295\times0.1067 = \qty{31.48}{mb}$, $832\times0.0373 =
\qty{31.03}{mb}$, $155\times0.1984 = \qty{30.75}{mb}$. (b)~Vanuit
$^{122}\mathrm{Te}$: $N_{123} = 31.48/832 = \num{0.03783}$, \qty{1.4}{\percent}
boven de gemeten waarde, en $N_{124} = 31.48/155 = \num{0.2031}$,
\qty{2.4}{\percent} boven. Abundanties die met een factor vijf verschillen, worden
binnen enkele procenten voorspeld uit de werkzame doorsneden alleen.
(c)~$^{128}\mathrm{Te}$ is alleen door het \omterm{def:b3:origin-of-elements:s-process}{r-proces} gevormd: de s-stroom eindigt bij tellurium in
$^{127}\mathrm{Te}$, dat vervalt, dus er gaat geen stationaire s-stroom door
$^{128}\mathrm{Te}$ en de abundantie werd bepaald door het \omterm{def:b3:origin-of-elements:s-process}{r-proces}.
\end{solution}

\begin{solution}{exo:b3:origin-of-elements:6}
(a)~$\tau_0\langle\sigma\rangle = 0.306\times155 = 47.4$, dus
$f = 1/(1 + 1/47.4) = \num{0.979}$ voor $^{124}\mathrm{Te}$; voor
$^{138}\mathrm{Ba}$ geldt $0.306\times4.0 = 1.224$ en $f = 1/(1 + 0.817) =
\num{0.550}$: de magische kern laat slechts \qty{55}{\percent} van de
stroom door. (b)~De zes factoren zijn \num{0.959}, \num{0.550}, \num{0.908},
\num{0.771}, \num{0.972} en \num{0.915}; hun product is \num{0.328}:
$\sigma N$ daalt met een factor \num{3.05}. (c)~De gemeten daling is
$27.6/7.72 = 3.57$. Het model geeft de grootte van de sprong correct binnen
\qty{15}{\percent}, bijna geheel afkomstig van $^{138}\mathrm{Ba}$ en
$^{140}\mathrm{Ce}$, de twee kernen met $N = 82$.
\end{solution}

\begin{solution}{exo:b3:origin-of-elements:7}
(a)~$kT = 8\times10^3\times\qty{1.602e-19}{J} = \qty{1.282e-15}{J}$, en
$v_T = \sqrt{2\times\num{1.282e-15}/\num{1.6749e-27}} =
\qty{1.24e6}{m.s^{-1}}$. (b)~Met $\langle\sigma\rangle =
\qty{155e-31}{m^2}$, $\lambda = n_\mathrm{n}\langle\sigma\rangle v_T$
bedraagt \qty{1.92e-11}{s^{-1}} voor $n_\mathrm{n} = \qty{e12}{m^{-3}}$
($10^6$ per kubieke centimeter): een gemiddelde tijd van \qty{5.2e10}{s}, ongeveer
\num{1650} jaar; dan \num{16.5} jaar bij $10^8$ en \num{0.165}~jaar,
twee maanden, bij \qty{e10}{cm^{-3}}. (c)~$^{127}\mathrm{Te}$ vervalt in
uren, veel sneller dan elk van deze invangtijden: het vervalt altijd en
het pad splitst daar niet. $^{134}\mathrm{Cs}$ leeft ongeveer twee
jaar, tussen de invangtijden bij $10^8$ en \qty{e10}{cm^{-3}}: bij
de hoge dichtheden van de heliumflitsen wordt een groot deel ervan
ingevangen voordat het vervalt. Het is een vertakkingspunt, en de abundanties aan weerszijden
ervan meten de neutronendichtheid.
\end{solution}

\begin{solution}{exo:b3:origin-of-elements:8}
(a)~$10^{6.30 - 3.14} = 1445$ keer zeldzamer dan calcium en
$10^{4.97 - 3.14} = 68$ keer zeldzamer dan titanium. (b)~$10^{8.46 - 1.38}
= \num{1.2e7}$. (c)~Scandium vertoont het oneven--even-effect: kernen met
een even aantal protonen zijn steviger gebonden en worden opgebouwd door het
invangen van heliumkernen in de brandingsfasen van massieve sterren, die
ze bevoordelen. Beryllium vertoont de dip van lithium, beryllium en boor:
deze broze kernen worden vernietigd bij de temperaturen in het inwendige van
sterren, en heliumbranding gaat van helium rechtstreeks naar koolstof,
zonder ze te vormen.
\end{solution}

\begin{solution}{exo:b3:origin-of-elements:9}
(a)~Laat $\alpha$ de \omterm{def:b3:origin-of-elements:alpha}{alfa-elementen} zijn en $\mathrm{Fe}_\mathrm{cc}$,
$\mathrm{Fe}_\mathrm{Ia}$ het ijzer uit de twee soorten supernova. Op
het plateau is $\alpha/\mathrm{Fe}_\mathrm{cc} = 10^{P}$ maal
de zonneverhouding. Bij $[\mathrm{Fe/H}] = 0$, $\alpha/(\mathrm{Fe}_\mathrm{cc} +
\mathrm{Fe}_\mathrm{Ia})$ is gelijk aan de zonneverhouding. Door te delen volgt
$\mathrm{Fe}_\mathrm{cc}/(\mathrm{Fe}_\mathrm{cc} + \mathrm{Fe}_\mathrm{Ia})
= 10^{-P}$, en het Type~Ia-aandeel is $1 - 10^{-P} = 1 - 10^{-0.4} =
\num{0.602}$. (b)~$10^{-P} = 1/3$ geeft $P = \log_{10}3 = 0.48$.
(c)~De Type~Ia-supernova's beginnen na dezelfde vertraging in beide sterrenstelsels.
Toen ze in het dwergstelsel begonnen, had het gas er slechts
$[\mathrm{Fe/H}] = -2$, tegen $-1$ in de Melkweg: het had zijn gas
tien keer zo traag in sterren omgezet.
\end{solution}

\begin{solution}{exo:b3:origin-of-elements:10}
(a)~Voor zirconium,
\[
\frac{(X/\mathrm{Mg})_\mathrm{Earth}}{(X/\mathrm{Mg})_\mathrm{CI}} =
\frac{10.5/228000}{3.96/95870} = 1.11 ;
\]
op dezelfde manier $(2670/228000)/(5010/95870) = 0.224$ voor natrium,
\num{0.190} voor kalium en \num{0.118} voor rubidium. (b)~Zirconium, het meest refractaire element, zit
dicht bij zijn chondritische verhouding. De drie matig vluchtige
alkalimetalen, die condenseren bij \qtyrange{800}{1006}{K}, zijn uitgeput
met factoren van vijf tot negen, rubidium, het meest vluchtige, het sterkst. De
silicaatrijke Aarde volgt de \omterm{def:b3:origin-of-elements:volatility}{vluchtigheidstrend}. (c)~Magnesium is een belangrijk,
refractair element dat met het grootste deel van het gesteente condenseert, en het is
lithofiel: het bleef in de mantel achter toen de kern zich vormde. Een siderofiele
referentie zoals ijzer of nikkel zou het verlies naar de kern vermengen met het
verlies door vluchtigheid.
\end{solution}

\begin{solution}{exo:b3:origin-of-elements:11}
(a)~Wanneer een massa $\mathrm{d}s$ in sterren wordt vastgelegd, verwijdert de wind
$\lambda\,\mathrm{d}s$ gas met \omterm{def:b3:origin-of-elements:metallicity}{metalliciteit} $Z$:
$\mathrm{d}M_\mathrm{g} = -(1 + \lambda)\,\mathrm{d}s$ en
$\mathrm{d}(ZM_\mathrm{g}) = y\,\mathrm{d}s - Z\,\mathrm{d}s -
\lambda Z\,\mathrm{d}s$. Na het uitwerken van de tweede vergelijking en het aftrekken van
$Z\,\mathrm{d}M_\mathrm{g}$ blijft $M_\mathrm{g}\,\mathrm{d}Z =
y\,\mathrm{d}s = -\frac{y}{1 + \lambda}\,\mathrm{d}M_\mathrm{g}$, wat
integreert tot het resultaat. (b)~$\ln(1/0.29) = 1.238$: $Z = 0.0155$ voor
$\lambda = 0$, en de helft daarvan, $0.0077$, voor $\lambda = 1$. (c)~Een wind
verlaagt de effectieve opbrengst met de factor $1 + \lambda$. In een sterrenstelsel
van kleine massa drijven supernova's het gas makkelijker naar buiten, $\lambda$ is groot,
en de metalen gaan verloren met het gas.
\end{solution}

\begin{solution}{exo:b3:origin-of-elements:12}
(a)~$N_\mathrm{Li}/N_\mathrm{H} = 10^{0.96 - 12} = \num{9.1e-12}$ in de
Zon en $10^{3.25 - 12} = \num{1.78e-9}$ in de CI-chondrieten. De Zon
houdt $10^{-2.29} = \num{0.0051}$ van zijn lithium over: hij heeft
\qty{99.5}{\percent} ervan vernietigd. (b)~$\num{1.78e-9}/\num{1.45e-10} = 12.3$ en
$\num{1.78e-9}/\num{4.72e-10} = 3.8$. (c)~Het Zonnestelsel bevat bijna
vier keer meer lithium dan de oerknal maakte: het meeste van zijn lithium
is later gemaakt, door een ander proces. De oudste sterren bevatten drie keer
minder dan de oerknal maakte: ofwel hebben ze een deel ervan vernietigd, zoals
de Zon heeft gedaan, ofwel mist de voorspelling iets; de twee vergelijkingen
alleen zeggen niet welke van beide.
\end{solution}

\begin{solution}{pb:b3:origin-of-elements:1}

\textbf{Deel I.}

\textbf{1.} $10^{12 - 7.46} = 10^{4.54} = \num{34700}$.

\textbf{2.} $N_\mathrm{He}/N_\mathrm{H} = 10^{10.914 - 12} = \num{0.0820}$, dus
$4N_\mathrm{He}/N_\mathrm{H} = \num{0.328}$, tegenover $Y/X =
0.2423/0.7438 = \num{0.326}$. Ze komen overeen tot op \qty{0.7}{\percent}; het
kleine verschil is er omdat een heliumatoom iets minder weegt dan vier
waterstofatomen.

\textbf{3.} $\log_{10}\bigl((1/350000)/\num{2.884e-5}\bigr) = -5.544 + 4.540 =
-1.00$.

\textbf{4.} $[\mathrm{O/H}] = -1.00 + 0.40 = -0.60$, dus
$\log\varepsilon(\mathrm{O}) = 8.69 - 0.60 = 8.09$. De ster heeft
$10^{8.09 - 6.46} = 42.7$ zuurstofatomen per ijzeratoom, tegenover
$10^{1.23} = 17.0$ in de Zon: $10^{0.4} = 2.5$ keer meer.

\textbf{5.} $Z \approx 0.0139\times10^{-0.60} = \num{0.0035}$, een kwart van de
Zon.

\textbf{6.} IJzer heeft heel veel lijnen in het spectrum van een koele ster, terwijl
zuurstof er weinig en zwakke heeft: ijzer is veel gemakkelijker en
nauwkeuriger te meten, en staat voor de metalen.

\textbf{Deel II.}

\textbf{7.} Elk ligt achter een stabiele isobaar van lagere $Z$:
$^{122}\mathrm{Sn}$, $^{123}\mathrm{Sb}$ en $^{124}\mathrm{Sn}$ stoppen
de vervalketens van het \omterm{def:b3:origin-of-elements:s-process}{r-proces} voordat ze telluur bereiken. Alle drie liggen
op het s-pad, dat via telluur loopt van $^{122}\mathrm{Te}$ tot
$^{126}\mathrm{Te}$.

\textbf{8.} \num{31.48}, \num{31.03} en \qty{30.75}{mb}; gemiddelde
\qty{31.09}{mb}.

\textbf{9.} Het gemiddelde van de andere twee is $(31.48 + 30.75)/2 = \qty{31.11}{mb}$,
dus $N_{123} = 31.11/832 = \num{0.03740}$, slechts \qty{0.3}{\percent}
boven de gemeten \num{0.0373}.

\textbf{10.} $1.329\times44.4 = \qty{59.0}{mb}$, bijna het dubbele van het plateau.
$^{128}\mathrm{Te}$ wordt alleen door het \omterm{def:b3:origin-of-elements:s-process}{r-proces} gevormd: het s-pad verlaat telluur bij
$^{127}\mathrm{Te}$, dus zijn overvloed is door het \omterm{def:b3:origin-of-elements:s-process}{r-proces} gemaakt en
heeft niets met de s-stroom te maken.

\textbf{11.} $1 - 7.43/8.53 = \num{0.13}$: \qty{13}{\percent} van de stroom ging
om $^{148}\mathrm{Sm}$ heen, door vangsten op de instabiele
$^{147}\mathrm{Nd}$, $^{147}\mathrm{Pm}$ en $^{148}\mathrm{Pm}$ voordat
ze tot samarium konden vervallen.

\textbf{12.} Met een gesloten schil van 82 neutronen binden deze kernen een extra
neutron zwak en hebben zeer kleine vangstdoorsneden (\qty{4.0}{mb}
voor $^{138}\mathrm{Ba}$). De stroom hoopt zich erop op en slechts een deel
gaat erdoorheen: het zijn knelpunten.

\textbf{Deel III.}

\textbf{13.} $\mu_1 = 13.7/(13.7 + 33.4) = 13.7/47.1 = \num{0.291}$.

\textbf{14.} $\ln(1/0.291) = \num{1.235}$, dus $y = 0.0154/1.235 = \num{0.0125}$,
\num{0.81} keer $Z_\odot$.

\textbf{15.} Bij $\mu = 0.9$ geldt $\ln(1/0.9) = 0.1054$, dus $Z = 0.0125\times0.1054 =
\num{0.00131}$. Omdat $0.00131/0.0154 = 0.0851$, geldt $[\mathrm{Fe/H}] =
-1.07$. Bij $\mu = 0.5$ geldt $Z = 0.0125\times0.693 = \num{0.00864}$ en
$[\mathrm{Fe/H}] = -0.25$.

\textbf{16.} $Z = 0.00154$, dus $\mu = \mathrm{e}^{-0.00154/0.01247} =
\mathrm{e}^{-0.1235} = \num{0.884}$.

\textbf{17.} $1 - 0.884 = \num{0.116}$: \qty{12}{\percent} van de massa.

\textbf{18.} $y = 0.0154/\ln 10 = \num{0.0067}$, ongeveer de helft.

\textbf{Deel IV.}

\textbf{19.} De massa van sterren met \omterm{def:b3:origin-of-elements:metallicity}{metalliciteit} onder $Z$ is
$M(1 - \mathrm{e}^{-Z/y})$, dus per eenheid $Z$ is zij
$(M/y)\,\mathrm{e}^{-Z/y}$. Omdat $Z = Z_\odot10^{[\mathrm{Fe/H}]}$ geldt,
$\mathrm{d}Z = Z\ln 10\,\mathrm{d}[\mathrm{Fe/H}]$, en per eenheid
$[\mathrm{Fe/H}]$ is de verdeling evenredig met
$Z\mathrm{e}^{-Z/y}$. De afgeleide,
$(1 - Z/y)\mathrm{e}^{-Z/y}$, verdwijnt bij $Z = y$: de piek ligt bij
$[\mathrm{Fe/H}] = \log_{10}0.810 = -0.09$.

\textbf{20.} $\num{0.116}/(1 - 0.291) = 0.1162/0.7091 = \num{0.164}$.

\textbf{21.} $Z/y = 1.235\times10^{-0.5} = 0.3905$, $1 - \mathrm{e}^{-0.3905} =
0.3233$, en $0.3233/0.7091 = \num{0.456}$.

\textbf{22.} $6670\times0.1638 = 1093$, tegenover 26: veertig keer te veel.

\textbf{23.} $26/902 = \qty{2.9}{\percent}$, nog steeds meer dan vijf keer lager dan
de doos. De sterren van het model zijn de langlevende, die van hooguit
ongeveer één zonsmassa, die allemaal sinds hun vorming nog bestaan; van
de massievere F-sterren zijn de oudste, die ook het metaalarmst zijn,
al dood, en de volledige steekproef telt ze te weinig.

\textbf{24.} $(1 - 0.1^{0.1})/(1 - 0.1) = 0.2057/0.9 = \num{0.23}$: slechter.
Een hogere gasfractie helpt evenmin: naarmate $\mu_1$ naar 1 gaat, nadert de
fractie $0.1$, nog steeds meer dan drie keer de
\qty{2.9}{\percent} die is waargenomen. Geen gasfractie redt de \omterm{thm:b3:origin-of-elements:closed-box}{gesloten doos}.

\textbf{25.} De \omterm{thm:b3:origin-of-elements:closed-box}{gesloten doos} voorspelt dat
\textbf{\qty[text-series-to-math=true]{16}{\percent}} van de langlevende
sterren in de zonnenabuurschap $[\mathrm{Fe/H}] < -1$ heeft, tegenover
hooguit \qty{3}{\percent} die is waargenomen. De schijf was geen \omterm{thm:b3:origin-of-elements:closed-box}{gesloten doos}:
zij groeide door een langzame instroom van metaalarm gas, zodat er zelden
gas lang bleef wachten op lage \omterm{def:b3:origin-of-elements:metallicity}{metalliciteit} om sterren te vormen.
\end{solution}
